% =====================================================================
\chapter{Guía para construir un optimizador con IA en cualquier ámbito}
\label{tema:6}
% =====================================================================

Este tema reúne la fase~VI en un procedimiento. No es un algoritmo, sino una disciplina de ingeniería: el orden en que conviene tomar las decisiones y los documentos que deben quedar escritos. Sigue la estructura del curso 16.842 «Fundamentals of Systems Engineering» del MIT~\cite{mit16842}, que recorre el modelo en V (partes interesadas y concepto de operación, requisitos con márgenes, arquitectura, exploración del espacio de alternativas, revisión preliminar del diseño), las clases~1 y~2 del curso CS329S de Stanford~\cite{cs329s} y el capítulo~2 de Huyen~\cite{huyen} para el planteamiento de sistemas con aprendizaje automático, y la lista de pruebas de Breck \emph{et al.}~\cite{breck}. Cada paso remite al módulo donde se estudió su contenido.

\section{Los doce pasos}

\begin{enumerate}[label=\textbf{\arabic*.},leftmargin=1.6em]
  \item \textbf{La decisión y sus responsables.} ¿Qué se decide, quién lo decide hoy, con qué frecuencia y con qué horizonte? ¿Quién usará el sistema, quién responde de sus errores y quién puede pararlo? Entregable: concepto de operación de una página.
  \item \textbf{Objetivo, indicadores y restricciones.} Un objetivo escalar, los indicadores que se vigilan sin optimizarlos, las restricciones duras (seguridad, ley, física) y las blandas, cada una con su unidad y su fuente (módulo~24, tema~1). Entregable: requisitos numéricos con márgenes.
  \item \textbf{Línea base y techo de valor.} Coste de la práctica actual, de la mejor alternativa sencilla sin IA y de una política con información perfecta (\cref{prop:techo}). Si el techo no compensa el coste del proyecto, se para aquí.
  \item \textbf{Datos.} Inventario de fuentes, frecuencias, calidad, retrasos, disponibilidad en el instante de decidir, propiedad y protección; versionado desde el primer día (\cref{tema:4}).
  \item \textbf{Formulación y familia de métodos.} El mapa del módulo~24 (tema~1, §1.3): programación matemática, optimización bayesiana para experimentos caros (módulo~25), control y RL para decisiones secuenciales (módulo~26), optimización geométrica (módulo~27), búsqueda de programas cuando lo que se busca es una regla (\cref{tema:1}).
  \item \textbf{Modelos predictivos o sustitutos.} Qué se predice, con qué horizonte y con qué métrica de decisión; validación temporal y selección con presupuesto (\cref{tema:3}).
  \item \textbf{Resolvedor, salvaguardas y respaldo.} Tiempo máximo de cálculo, comprobación independiente de las restricciones duras, política de respaldo.
  \item \textbf{Evaluación fuera de línea.} Sobre históricos (sin fugas), en un gemelo con números aleatorios comunes y, si hay registros con propensiones, con estimadores contrafactuales (\cref{tema:4}).
  \item \textbf{Piloto.} Modo sombra y después experimento con tamaño calculado (\cref{prop:tamano}) y criterio de aceptación escrito antes de empezar.
  \item \textbf{Despliegue.} API con contrato, interfaz que explica y permite anular, registro de decisiones, versiones y vuelta atrás.
  \item \textbf{Monitorización y reentrenamiento.} Entradas, predicciones, decisiones, anulaciones e indicadores; disparadores y periodo de reentrenamiento (\cref{prop:reentreno}).
  \item \textbf{Valor, riesgo, normativa y ciclo de vida.} Intervalos para la mejora, VAN, AMFE, encaje regulatorio, documentación y criterios de retirada (\cref{tema:5}).
\end{enumerate}

\begin{proposicion}[Techo de valor]\label{prop:techo}
Sea $f(x,\theta)$ el coste de la decisión $x\in\mathcal{X}$ cuando el estado incierto es $\theta$, y sea $\hat x$ cualquier regla de decisión que solo use información disponible antes de decidir. Entonces
\[
\E\bigl[f(\hat x,\theta)\bigr]\ \ge\ \E\Bigl[\min_{x\in\mathcal{X}}f(x,\theta)\Bigr].
\]
En consecuencia, ninguna mejora de la predicción puede reducir el coste esperado de la mejor regla sin información por debajo del coste con información perfecta, y la diferencia entre ambos (el valor de la información perfecta del módulo~24, tema~6, definición~6.6) acota lo que puede ganarse prediciendo mejor.
\end{proposicion}

En el caso del mantenimiento, una política que conoce el futuro e interviene cada máquina en la última ventana antes de que fuera a fallar cuesta de media $76\,170$~€/año en el gemelo (sin límite de técnicos), frente a $193\,670$ de la política por edad y $159\,300$ del optimizador. Este solo captura un 29\,\% del margen que deja la alternativa sencilla frente a la información perfecta: hay espacio para mejores sensores o mejores modelos, y el paso~3 permite estimar cuánto valdrían antes de pagarlos.

\section{Plantilla de documento de diseño}

El \cref{tab:plantilla} es la plantilla que se rellena antes de escribir código y se actualiza en cada revisión. Su objetivo no es la burocracia, sino forzar que las decisiones difíciles (qué es éxito, qué no se puede violar, contra qué se compara) se tomen pronto y por escrito, como en la revisión preliminar de diseño de 16.842~\cite{mit16842}.

\begin{table}[h]
\centering\small
\caption{Plantilla del documento de diseño de un optimizador con IA.}\label{tab:plantilla}
\begin{tabularx}{\textwidth}{@{}>{\raggedright\arraybackslash}p{3.6cm}>{\raggedright\arraybackslash}X@{}}
\toprule
Sección & Preguntas que debe responder \\
\midrule
1. Contexto y decisión & ¿Qué se decide, cada cuánto, con qué horizonte? ¿Qué pasa hoy? \\
2. Partes interesadas & ¿Quién decide, quién ejecuta, quién responde, quién puede parar el sistema? \\
3. Objetivo y requisitos & Objetivo, indicadores, restricciones duras y blandas con unidades; requisitos de latencia, disponibilidad y precisión con márgenes. \\
4. Línea base y techo & Coste actual, mejor alternativa sencilla, coste con información perfecta; valor máximo alcanzable. \\
5. Datos & Fuentes, calidad, retrasos, disponibilidad en el instante de decidir, versionado, protección. \\
6. Formulación y método & Modelo matemático; familia elegida y alternativas descartadas con su motivo. \\
7. Modelos y validación & Qué se predice; división temporal; métricas de predicción y de decisión; selección. \\
8. Salvaguardas y respaldo & Comprobaciones independientes; política de respaldo; tiempo máximo de cálculo. \\
9. Evaluación fuera de línea & Históricos, gemelo, contrafactual; resultados con intervalos. \\
10. Piloto & Diseño, tamaño, duración, criterio de aceptación fijado de antemano. \\
11. Despliegue y supervisión & API, interfaz, modos de operación, registro, vuelta atrás. \\
12. Monitorización & Señales, umbrales, acciones, reentrenamiento. \\
13. Valor y riesgo & VAN con incertidumbre, AMFE, pruebas de estrés, encaje normativo. \\
14. Ciclo de vida & Responsable del mantenimiento, revisiones periódicas, criterios de retirada. \\
\bottomrule
\end{tabularx}
\end{table}

\section{Mini-caso A: secuenciación de formatos en una fábrica de bebidas}

\addtocounter{definicion}{1}%

\section{Mini-caso B: presupuesto de geometría en un videojuego}

\addtocounter{definicion}{1}%

\section{Mini-caso C: asignación de trabajos en un centro de datos}

\addtocounter{definicion}{1}%

\begin{table}[h]
\centering\small
\caption{Tecnologías del tema~6: pila de referencia de un optimizador con IA.}\label{tab:tec6}
\begin{tabularx}{\textwidth}{@{}>{\raggedright\arraybackslash}p{2.5cm}>{\raggedright\arraybackslash}X>{\raggedright\arraybackslash}p{2.1cm}>{\raggedright\arraybackslash}X@{}}
\toprule
Capa & Herramientas habituales & Licencia o tipo & Cuándo elegirla \\
\midrule
Datos y versionado & DVC~\cite{dvc}, Great Expectations & abiertas (Apache~2.0) & desde el primer prototipo \\
Modelos & scikit-learn, PyTorch; Optuna~\cite{optuna} para el ajuste & abiertas (BSD-3; BSD-3; MIT) & según el tipo de datos (\cref{tema:3}) \\
Decisión & SciPy/HiGHS, OR-Tools, CVXPY (módulo~24); BoTorch (módulo~25) & abiertas & según la estructura del problema \\
Búsqueda con LLM & OpenEvolve~\cite{openevolve}, DSPy~\cite{dspy} & abiertas (Apache~2.0; MIT) & reglas o cadenas de llamadas con evaluador fiable \\
Servicio & FastAPI~\cite{fastapi}, contenedores, Kubernetes~\cite{k8s} & abiertas (MIT; Apache~2.0) & integración con otros sistemas \\
Seguimiento & MLflow~\cite{mlflow}, Evidently~\cite{evidently} & abiertas (Apache~2.0) & trazabilidad y deriva \\
Escala & Ray~\cite{ray} & abierta (Apache~2.0) & simulaciones y ajustes en paralelo \\
\bottomrule
\end{tabularx}
\end{table}

\addtocounter{definicion}{1}%

